\begin{table*}[t]
\centering
\caption{Vollständige Kontraste der separaten RegNet- und vollständigen ResNet-Experimente mit jeweils vierzig Prozessen. Geometrische Mittel haben nominale 95\%-Intervalle über fünf vollständige Blöcke. SD bezeichnet die Stichprobenstandardabweichung der logarithmierten Blockkontraste. Die Intervalle haben vier Freiheitsgrade.}
\label{tab:supp-scheduling-summary}
\begingroup\small\setlength{\tabcolsep}{3pt}
\begin{tabular}{lllll}
\toprule
Kontrast & RegNet-Mittel [Intervall] & Log-SD & ResNet-Mittel [Intervall] & Log-SD \\
\midrule
Generic produktiv & 0.6603 [0.6559, 0.6648] & 0.00541 & 0.6057 [0.6033, 0.6082] & 0.00328 \\
Native produktiv & 3.3713 [3.3555, 3.3872] & 0.00378 & 2.0597 [2.0281, 2.0919] & 0.01248 \\
Generic Einzelauftrag & 1.0289 [1.0228, 1.0350] & 0.00478 & 0.8359 [0.8276, 0.8442] & 0.00804 \\
Native Einzelauftrag & 1.0283 [1.0163, 1.0405] & 0.00947 & 0.8290 [0.8240, 0.8339] & 0.00481 \\
Generic Scheduling & 0.6418 [0.6350, 0.6487] & 0.00856 & 0.7247 [0.7183, 0.7310] & 0.00708 \\
Native Scheduling & 3.2785 [3.2295, 3.3282] & 0.01213 & 2.4847 [2.4382, 2.5320] & 0.01520 \\
Interaktion & 5.1081 [4.9896, 5.2294] & 0.01890 & 3.4288 [3.3532, 3.5061] & 0.01796 \\
\bottomrule
\end{tabular}
\endgroup
\end{table*}
